\chapter{Eye Twitching (Blepharospasm and Myokymia)}

\section*{Definition}
Intermittent, involuntary, repetitive contractions of the orbicularis oculi muscle. Benign eyelid myokymia involves fine, rippling twitching of a single eyelid (usually lower), while blepharospasm involves sustained, bilateral, forceful eyelid closure.

\section*{What Happens in Your Body}
Eyelid myokymia results from spontaneous depolarization of motor units in the orbicularis oculi muscle, often triggered by fatigue, stress, or caffeine. The pathophysiology involves increased excitability of the facial nerve nucleus or motor cortex. Blepharospasm is a focal dystonia involving involuntary contraction of the orbicularis oculi due to basal ganglia dysfunction and altered inhibition in the brainstem blink reflex circuits.

\section*{Causes}
\begin{itemize}
  \item Fatigue and sleep deprivation (most common cause of myokymia)
  \item Stress and anxiety
  \item Caffeine or stimulant intake
  \item Eye strain (prolonged screen time, uncorrected vision)
  \item Dry eyes (ocular surface irritation triggers reflex blinking)
  \item Nutritional deficiencies (magnesium, potassium, calcium)
  \item Alcohol or tobacco use
  \item Benign essential blepharospasm (focal dystonia, typically bilateral)
  \item Hemifacial spasm (unilateral, involving entire side of face, vascular compression of CN VII)
  \item Medications (antipsychotics, antihistamines, SSRIs)
\end{itemize}

\section*{When to See a Doctor}
See your doctor if:
\begin{itemize}
  \item Symptoms are severe or getting worse over time
  \item Twitching spreading to other parts of the face, both eyes closing involuntarily, eye redness or discharge, or twitching lasting more than 2 weeks
  \item You have other concerning symptoms such as fever, weight loss, or bleeding
\end{itemize}

\section*{Related Symptoms}
\begin{itemize}
  \item Eye strain
  \item Dry eyes
  \item Facial twitching
  \item Blepharospasm
  \item Fatigue
\end{itemize}
\textbf{Important:} If this symptom persists, your doctor will check for serious underlying causes.

\section*{Self-Care / Home Management}
\begin{itemize}
  \item Get adequate sleep and reduce fatigue, the most common trigger for eyelid myokymia.
  \item Reduce caffeine intake from coffee, tea, soda, and chocolate.
  \item Apply warm compresses to the eyelid and gently massage the area.
  \item Use artificial tears if dry eye is contributing to reflex twitching.
  \item Practice stress reduction techniques such as meditation, deep breathing, or yoga.
\end{itemize}

\section*{How Your Doctor Will Evaluate You}
\begin{itemize}
  \item Your doctor will take a history of the twitching pattern, triggers, and associated symptoms.
  \item A neurological examination assesses for hemifacial spasm or other cranial nerve abnormalities.
  \item Slit-lamp examination checks for dry eye, blepharitis, or ocular surface irritation.
  \item Blood tests may check electrolyte levels (magnesium, potassium, calcium).
  \item MRI with angiography may be ordered if hemifacial spasm or structural compression is suspected.
\end{itemize}

\section*{Treatment Options}
\begin{itemize}
  \item Reassurance and trigger avoidance (sleep, caffeine, stress) for benign myokymia.
  \item Botulinum toxin injections are first-line treatment for blepharospasm and hemifacial spasm.
  \item Treatment of dry eye with artificial tears and lid hygiene if contributing.
  \item Electrolyte supplementation if deficiencies are identified.
  \item Microvascular decompression surgery for severe hemifacial spasm from vascular compression.
\end{itemize}


\section*{References}
\begin{thebibliography}{99}
\bibitem{jankovic2016}
Jankovic J. Blepharospasm. In: \textit{Movement Disorders}. McGraw-Hill; 2016.
\bibitem{hallett2010}
Hallett M, Evinger C, Jankovic J, Stacy M. Update on blepharospasm. \textit{Neurology}. 2008;71(16):1275--1282.
\bibitem{simpson2015}
Simpson DM, Hallett M, Ashman EJ, et al. Practice guideline update: botulinum neurotoxin for the treatment of blepharospasm, cervical dystonia, and other movement disorders. \textit{Neurology}. 2016;86(19):1818--1826.
\bibitem{defazio2017}
Defazio G, Hallett M, Jinnah HA, Conte A, Berardelli A. Blepharospasm and other cranial dystonias. \textit{Handb Clin Neurol}. 2017;141:525--537.
\bibitem{opdycke2007}
Opdycke R. Eyelid twitching. In: \textit{Clinical Neurology for Psychiatrists}. 6th ed. Saunders; 2007.
\bibitem{benito2019}
Benito-Le\'{o}n J, Louis ED. Essential tremor and benign eyelid twitching. \textit{Tremor Other Hyperkinet Mov}. 2019;9.
\end{thebibliography}
